% ============================================================
\chapter{Introduction to MLOps --- From Notebook to Production}
\label{ch:introduction}
\index{MLOps!introduction}
% ============================================================

\section{What is MLOps?}
\label{sec:mlops-definition}

\textbf{MLOps}\index{MLOps} refers to the set of practices, tools, and
cultural principles aimed at deploying and maintaining machine learning
models in production reliably and efficiently. The term is a portmanteau
of \emph{Machine Learning} and \emph{Operations}, by analogy with
\emph{DevOps}\index{DevOps} from software engineering.

\begin{definition}[MLOps]
MLOps is the discipline that applies DevOps principles (continuous
integration, continuous deployment, monitoring, automation) to the lifecycle
of machine learning models, adding specificities related to data,
experiments, and model drift.
\end{definition}

\subsection{DevOps vs MLOps}

Traditional DevOps manages code. MLOps manages code \textbf{and} data
\textbf{and} models---three artifacts whose interaction creates additional
complexity.

\begin{center}
\begin{tikzpicture}[
  box/.style={rectangle, draw, rounded corners, minimum width=3cm,
    minimum height=1.5cm, align=center, font=\small},
  arr/.style={-{Stealth[length=2.5mm]}, thick}
]
  \node[box, fill=blue!10] (devops) {\textbf{DevOps}\\Code};
  \node[box, fill=green!10, right=3cm of devops] (mlops) {\textbf{MLOps}\\Code + Data\\+ Models};

  \node[above=1.5cm of devops, font=\small\bfseries] (shared) {Shared Practices};
  \node[below=0.3cm of shared, font=\scriptsize, align=center] {CI/CD, containerization,\\monitoring, IaC};

  \draw[arr, steelblue] (shared.south) -- (devops.north);
  \draw[arr, steelblue] (shared.south) -- (mlops.north);

  \node[below=0.8cm of mlops, font=\scriptsize, align=center, text=deepred] {+ Data versioning\\+ Experiment tracking\\+ Drift detection\\+ Feature engineering};
\end{tikzpicture}
\end{center}

\begin{remark}
The fundamental difference is that the behaviour of an ML system depends
not only on the code but also on the data. The same code trained on
different data produces a different model. This requires versioning data
just like code.
\end{remark}

\section{MLOps Maturity Levels}
\label{sec:maturity-levels}
\index{MLOps!maturity levels}

Google proposed three maturity levels for ML systems in production:

\begin{center}
\begin{tabular}{@{}llp{7cm}@{}}
\toprule
\textbf{Level} & \textbf{Name} & \textbf{Description} \\
\midrule
0 & Manual & Entirely manual process. Notebooks, no pipeline. Complete
  separation between data scientists and ops. \\
1 & ML Pipeline & Automated training pipeline. Continuous model deployment.
  Basic monitoring. \\
2 & CI/CD Pipeline & CI/CD pipeline for ML code. Automated training
  triggered by data. Automatic retraining on drift. \\
\bottomrule
\end{tabular}
\end{center}

\begin{center}
\begin{tikzpicture}[
  level/.style={rectangle, draw, rounded corners, minimum width=10cm,
    minimum height=1.2cm, align=center, font=\small},
  arr/.style={-{Stealth[length=2.5mm]}, thick, steelblue}
]
  \node[level, fill=red!10] (l0) {\textbf{Level 0: Manual}\\
    Notebooks $\to$ export model $\to$ manual deployment};
  \node[level, fill=yellow!10, above=0.5cm of l0] (l1) {\textbf{Level 1: ML Pipeline}\\
    Automated pipeline $\to$ continuous deployment $\to$ monitoring};
  \node[level, fill=green!10, above=0.5cm of l1] (l2) {\textbf{Level 2: CI/CD + ML}\\
    Code + data CI/CD $\to$ auto retraining $\to$ A/B testing};

  \draw[arr] (l0) -- (l1);
  \draw[arr] (l1) -- (l2);
\end{tikzpicture}
\end{center}

\section{The ML Lifecycle}
\label{sec:ml-lifecycle}
\index{ML lifecycle}

An ML project in production goes through several phases that form a cycle
(not a linear sequence):

\begin{enumerate}
  \item \textbf{Problem definition}: what is the business objective? What
    is the success criterion?
  \item \textbf{Data collection and preparation}: acquisition, cleaning,
    labelling, feature engineering.
  \item \textbf{Experimentation}: model selection, hyperparameter
    optimization, evaluation.
  \item \textbf{Pipeline development}: transforming the notebook into
    modular, tested, versioned code.
  \item \textbf{Deployment}: putting the model into production (API,
    batch, edge).
  \item \textbf{Monitoring}: tracking performance, detecting drift, alerts.
  \item \textbf{Retraining}: updating the model when data or performance
    evolve.
\end{enumerate}

\begin{center}
\begin{tikzpicture}[
  node distance=1.8cm,
  phase/.style={circle, draw, minimum size=1.6cm, align=center,
    font=\scriptsize\bfseries, inner sep=1pt},
  arr/.style={-{Stealth[length=2mm]}, thick, steelblue}
]
  \node[phase, fill=blue!15] (prob) {Problem};
  \node[phase, fill=blue!10, right=1.5cm of prob] (data) {Data};
  \node[phase, fill=green!10, below right=1cm and 1cm of data] (exp) {Experim.};
  \node[phase, fill=green!15, below=1.5cm of exp] (pipe) {Pipeline};
  \node[phase, fill=orange!10, below left=1cm and 1cm of pipe] (deploy) {Deploy};
  \node[phase, fill=red!10, left=1.5cm of deploy] (monitor) {Monitor};
  \node[phase, fill=red!15, above left=1cm and 1cm of monitor] (retrain) {Retrain};

  \draw[arr] (prob) -- (data);
  \draw[arr] (data) -- (exp);
  \draw[arr] (exp) -- (pipe);
  \draw[arr] (pipe) -- (deploy);
  \draw[arr] (deploy) -- (monitor);
  \draw[arr] (monitor) -- (retrain);
  \draw[arr] (retrain) -- (data);
\end{tikzpicture}
\end{center}

\section{Technical Debt in ML}
\label{sec:technical-debt}
\index{technical debt}

Sculley et al.~\cite{sculley2015} identified that ML code represents only a
small fraction of a production ML system. The rest is infrastructure:

\begin{center}
\begin{tikzpicture}[
  box/.style={rectangle, draw, minimum height=0.8cm, align=center,
    font=\scriptsize}
]
  \node[box, fill=black!80, text=white, minimum width=2cm] (ml) {ML Code};
  \node[box, fill=yellow!20, minimum width=3.5cm, above=0.1cm of ml] (config) {Configuration};
  \node[box, fill=yellow!20, minimum width=3.5cm, below=0.1cm of ml] (data) {Data Collection};
  \node[box, fill=yellow!20, minimum width=1.5cm, left=0.1cm of ml] (feat) {Feature\\Engineering};
  \node[box, fill=yellow!20, minimum width=1.5cm, right=0.1cm of ml] (serve) {Serving\\Infrastructure};
  \node[box, fill=orange!20, minimum width=6.5cm, above=0.1cm of config] (monitor) {Monitoring \quad Process Tools \quad Resource Management};
  \node[box, fill=orange!20, minimum width=6.5cm, below=0.1cm of data] (valid) {Data Verification \quad Analysis \quad Testing};
\end{tikzpicture}
\end{center}

\begin{warning}[title=\textbf{The notebook trap}]
Jupyter notebooks are excellent for exploration but poor for production.
Common problems:
\begin{itemize}
  \item Non-linear execution (cells run out of order)
  \item Hidden state (global variables that persist)
  \item No unit tests
  \item Difficult to version (large JSON files)
  \item No modularity (functions and classes not reusable)
\end{itemize}
\end{warning}

\section{MLOps System Architecture}
\label{sec:mlops-architecture}

A mature MLOps system includes the following components:

\begin{toolbox}[title=\textbf{MLOps system components}]
\begin{description}
  \item[Code management] Git, GitHub/GitLab
  \item[Data management] DVC, Delta Lake, LakeFS
  \item[Environment management] Conda, pip, venv, Poetry
  \item[Experiment tracking] MLflow, Weights \& Biases, Neptune
  \item[Training pipeline] Airflow, Prefect, Kubeflow Pipelines
  \item[Feature store] Feast, Tecton, Hopsworks
  \item[Containerization] Docker, Kubernetes
  \item[Serving] FastAPI, TorchServe, TensorFlow Serving, Triton
  \item[CI/CD] GitHub Actions, GitLab CI, Jenkins
  \item[Monitoring] Prometheus, Grafana, Evidently, WhyLabs
\end{description}
\end{toolbox}

\section{Tutorial: From Notebook to Modular Script}
\label{sec:notebook-to-script}

\subsection{The starting notebook}

Consider a typical classification notebook:

\begin{codeblock}[title=\textbf{Typical notebook (anti-pattern)}]
\begin{minted}{python}
# Cell 1
import pandas as pd
from sklearn.ensemble import RandomForestClassifier
from sklearn.model_selection import train_test_split
from sklearn.metrics import accuracy_score

# Cell 2
df = pd.read_csv("data.csv")
X = df.drop("target", axis=1)
y = df["target"]

# Cell 3
X_train, X_test, y_train, y_test = train_test_split(
    X, y, test_size=0.2, random_state=42
)

# Cell 4
model = RandomForestClassifier(n_estimators=100, random_state=42)
model.fit(X_train, y_train)
print(f"Accuracy: {accuracy_score(y_test, model.predict(X_test)):.4f}")
\end{minted}
\end{codeblock}

\subsection{The modular script}

\begin{bestpractice}[title=\textbf{Recommended project structure}]
\begin{minted}{text}
my_ml_project/
  src/
    __init__.py
    data.py          # Loading and preprocessing
    model.py         # Model definition
    train.py         # Training loop
    evaluate.py      # Evaluation metrics
  tests/
    test_data.py
    test_model.py
  configs/
    config.yaml      # Hyperparameters
  requirements.txt
  Makefile
  README.md
\end{minted}
\end{bestpractice}

\begin{codeblock}[title=\textbf{src/data.py --- Data module}]
\begin{minted}{python}
"""Data loading and preprocessing module."""
import pandas as pd
from sklearn.model_selection import train_test_split


def load_data(path: str) -> pd.DataFrame:
    """Load dataset from CSV file."""
    df = pd.read_csv(path)
    return df


def split_data(
    df: pd.DataFrame,
    target_col: str = "target",
    test_size: float = 0.2,
    random_state: int = 42,
) -> tuple:
    """Split data into train and test sets."""
    X = df.drop(target_col, axis=1)
    y = df[target_col]
    return train_test_split(
        X, y, test_size=test_size, random_state=random_state
    )
\end{minted}
\end{codeblock}

\begin{codeblock}[title=\textbf{src/train.py --- Training script}]
\begin{minted}{python}
"""Training script with configuration."""
import argparse
import yaml
import joblib
from data import load_data, split_data
from sklearn.ensemble import RandomForestClassifier
from sklearn.metrics import accuracy_score


def train(config: dict) -> None:
    """Train model with given configuration."""
    df = load_data(config["data"]["path"])
    X_train, X_test, y_train, y_test = split_data(
        df,
        target_col=config["data"]["target_col"],
        test_size=config["data"]["test_size"],
    )

    model = RandomForestClassifier(**config["model"])
    model.fit(X_train, y_train)

    y_pred = model.predict(X_test)
    acc = accuracy_score(y_test, y_pred)
    print(f"Accuracy: {acc:.4f}")

    joblib.dump(model, config["output"]["model_path"])
    print(f"Model saved to {config['output']['model_path']}")


if __name__ == "__main__":
    parser = argparse.ArgumentParser()
    parser.add_argument("--config", default="configs/config.yaml")
    args = parser.parse_args()

    with open(args.config) as f:
        config = yaml.safe_load(f)

    train(config)
\end{minted}
\end{codeblock}

\begin{codeblock}[title=\textbf{configs/config.yaml}]
\begin{minted}{yaml}
data:
  path: "data/dataset.csv"
  target_col: "target"
  test_size: 0.2

model:
  n_estimators: 100
  max_depth: 10
  random_state: 42

output:
  model_path: "models/rf_model.joblib"
\end{minted}
\end{codeblock}

\begin{codeblock}[title=\textbf{Makefile}]
\begin{minted}{makefile}
.PHONY: install train test clean

install:
	pip install -r requirements.txt

train:
	python src/train.py --config configs/config.yaml

test:
	pytest tests/ -v

clean:
	rm -rf models/*.joblib __pycache__ .pytest_cache
\end{minted}
\end{codeblock}

\section{Best Practices and Anti-patterns}

\begin{bestpractice}[title=\textbf{Fundamental principles}]
\begin{enumerate}
  \item \textbf{Separation of concerns}: data, model, training, evaluation
    in separate modules.
  \item \textbf{Externalized configuration}: hyperparameters in YAML files,
    not hardcoded.
  \item \textbf{Reproducibility by default}: fix random seeds, version data
    and code.
  \item \textbf{Automated tests}: test data loading, tensor shapes,
    convergence on a small sample.
  \item \textbf{Automation}: Makefile or scripts for all repetitive
    operations.
\end{enumerate}
\end{bestpractice}

\begin{antipattern}[title=\textbf{Absolutely avoid}]
\begin{itemize}
  \item Hardcoding absolute paths: \cli{/home/alice/data/train.csv}
  \item Committing large datasets to Git
  \item Committing secrets (API keys, passwords)
  \item Working exclusively in notebooks without modular code
  \item Ignoring tests because ``it's ML''
  \item Deploying a model without monitoring
\end{itemize}
\end{antipattern}

\section{Mini-project: Structuring an ML Project}
\label{sec:mini-project-ch1}

\begin{miniproject}[title=\textbf{Mini-project 1: From notebook to structured project}]
\begin{enumerate}
  \item Create the recommended directory structure.
  \item Transform the classification notebook above into separate Python
    modules.
  \item Write a \file{config.yaml} file for hyperparameters.
  \item Create a \file{Makefile} with targets \cli{install}, \cli{train},
    \cli{test}, \cli{clean}.
  \item Write at least two unit tests (one for data loading, one for model
    output shape).
  \item Verify that \cli{make train} works end-to-end.
\end{enumerate}
\end{miniproject}

\section{Exercises}
\label{sec:exercises-ch1}

\begin{exercise}[$\bigstar$ --- Project structure]
Create an ML project skeleton for a regression problem with the recommended
structure. Include a \file{README.md} describing the project and a
\file{.gitignore} adapted for ML (ignore \file{data/}, \file{models/},
\file{*.pyc}, \file{.ipynb\_checkpoints/}).
\end{exercise}

\begin{exercise}[$\bigstar\bigstar$ --- Notebook refactoring]
Take an existing Jupyter notebook (your own or a Kaggle notebook) and
transform it into a structured project:
\begin{enumerate}
  \item Identify reusable functions
  \item Create the corresponding Python modules
  \item Externalize the configuration
  \item Write tests for critical functions
\end{enumerate}
\end{exercise}

\begin{exercise}[$\bigstar\bigstar\bigstar$ --- Complete project]
Build a complete ML project for text classification (e.g., sentiment
analysis on the IMDb dataset):
\begin{enumerate}
  \item Complete modular structure
  \item YAML configuration
  \item Unit tests with \pkg{pytest}
  \item Training script with command-line arguments
  \item Automatic metric reporting (precision, recall, F1)
  \item Complete \file{Makefile}
\end{enumerate}
\end{exercise}

\section{Cheatsheet}

\begin{cheatsheet}[title=\textbf{Chapter 1 Summary}]
\begin{tabular}{@{}p{4cm}p{8.5cm}@{}}
\toprule
\textbf{Concept} & \textbf{Key Point} \\
\midrule
MLOps & DevOps + data + models \\
Maturity levels & 0 (manual) $\to$ 1 (pipeline) $\to$ 2 (CI/CD) \\
Technical debt & ML code is a small fraction of the system \\
Notebook anti-pattern & Hidden state, no tests, no modularity \\
Project structure & \file{src/}, \file{tests/}, \file{configs/}, Makefile \\
External config & YAML for hyperparameters \\
\bottomrule
\end{tabular}
\end{cheatsheet}
